\documentclass[10pt,a4paper]{article}
\usepackage{amsmath,amssymb,amsthm,graphicx}
\usepackage{amsmath,amssymb,amsthm,amsfonts}
\usepackage{latexsym,bbm,graphicx,float,mathtools}
\usepackage{subcaption}

\usepackage{verbatim}
\usepackage{enumitem}
\usepackage{hyperref}  
\hypersetup{
    colorlinks,
    citecolor=black,
    filecolor=black,
    linkcolor=black,
    urlcolor=black,
    linktoc=all
}

\usepackage[noend]{algpseudocode}
\usepackage{algorithm}

\usepackage{thmtools}
\usepackage{thm-restate}

%primed-thm
\newtheorem{manualtheoreminner}{Theorem}
\newenvironment{manualtheorem}[1]{%
  \renewcommand\themanualtheoreminner{#1}%
  \manualtheoreminner
}{\endmanualtheoreminner}

%footnote twice preamble
%https://tex.stackexchange.com/questions/35043/reference-different-places-to-the-same-footnote
\makeatletter
\newcommand\footnoteref[1]{\protected@xdef\@thefnmark{\ref{#1}}\@footnotemark}
\makeatother



 % environments
\newtheorem{problem}{Problem}
\newtheorem{theorem}{Theorem}[section]
\newtheorem{lem}[theorem]{Lemma}

\newtheorem{proposition}[theorem]{Proposition}
\newtheorem{fact}[theorem]{Fact}
\newtheorem{cor}[theorem]{Corollary}
\newtheorem{conjecture}[theorem]{Conjecture}
\newtheorem{question}{Question}
\newtheorem{goal}{Goal}

% non-italicized environments
\theoremstyle{definition}%\newtheorem{definition}[theorem]{Definition}
%\newtheorem{remark}[theorem]{Remark}
\newtheorem{remark}[theorem]{Remark}
\newtheorem{claim}[theorem]{Claim}
\newtheorem{definition}[theorem]{Definition}
\newtheorem{obs}[theorem]{Observation}
\newtheorem{notation}[theorem]{Notation}
\newtheorem{example}[theorem]{Example}
\newtheorem{examples}[theorem]{Examples}
\theoremstyle{plain}

% probability
\newcommand{\std}{\operatorname{{\bf std}}}
\newcommand{\Var}{\operatorname{{\bf Var}}}
\newcommand{\Cov}{\operatorname{{\bf Cov}}}
\newcommand{\Varx}{\mathop{{\bf Var}\/}}
%\newcommand{\E}{\operatorname{{\mathbb E}}}
\DeclareMathOperator*{\E}{\mathbb{E}}
%\renewcommand{\Pr}{\operatorname{{\mathbb P}}}
\let\Pr\relax
\DeclareMathOperator*{\Pr}{\mathbb{P}}
\newcommand{\Bin}{\operatorname{Bin}}
\newcommand{\Ber}{\operatorname{Ber}}

% mathrm terms
\newcommand{\sbp}{\sum\mathrm{smABP}}
\newcommand{\poly}{\mathrm{poly}}
\newcommand{\imm}[2]{\mathrm{IMM}_{#1,#2}}
\newcommand{\rk}{\mathrm{relrk}}
\newcommand{\polylog}{\mathrm{polylog}}
\newcommand{\size}{\mathrm{size}}
\newcommand{\avg}{\mathop{\mathrm{avg}}}
\newcommand{\odd}{\mathrm{odd}}
\newcommand{\even}{\mathrm{even}}
\newcommand{\sign}{\mathrm{sign}}
\newcommand{\dist}{\mathrm{dist}}
\newcommand{\dtv}{d_\mathrm{TV}}
\newcommand{\supp}{\mathrm{supp}}
\newcommand{\Vol}{\mathrm{Vol}}
\newcommand{\opt}{\mathrm{opt}}
\newcommand{\rank}{\mathrm{rank}}
\newcommand{\Tr}{\mathrm{Tr}}
\newcommand{\Alg}{\mathrm{Alg}}
\newcommand{\sen}{\mathrm{sen}}
\newcommand{\dnf}{\mathrm{dnf}}
\newcommand{\dt}{\mathrm{dt}}
\newcommand{\type}{\mathrm{type}}
\newcommand{\cent}{\mathrm{center}}
\newcommand{\dyck}{\mathrm{Dyck}}
\newcommand{\out}{\mathrm{out}}
\newcommand{\children}{\mathrm{children}}
\newcommand{\Cop}{\mathsf{Copies}}

%mathsf complexity classes
\newcommand{\VP}{\mathsf{VP}}
\newcommand{\VBP}{\mathsf{VBP}}
\newcommand{\VNP}{\mathsf{VNP}}
\newcommand{\VF}{\mathsf{VF}}
%\newcommand{\P}{\mathsf{P}}
\newcommand{\NP}{\mathsf{NP}}




% mathbf terms
\newcommand{\Size}{\mathbf{Size}}
\newcommand{\Depth}{\mathbf{Depth}}

% number systems
\newcommand{\R}{\mathbb R}
\newcommand{\C}{\mathbb C}
\newcommand{\N}{\mathbb N}
\newcommand{\Z}{\mathbb Z}
\newcommand{\F}{\mathbb F}
\newcommand{\K}{\mathbb K}
\newcommand{\Q}{\mathbb Q}

% short forms
\newcommand{\eps}{\varepsilon}
\renewcommand{\epsilon}{\eps} 
\newcommand{\la}{\langle}
\newcommand{\ra}{\rangle}
\newcommand{\wt}[1]{\widetilde{#1}}
\newcommand{\wh}[1]{\widehat{#1}}
\newcommand{\sse}{\subseteq}
\newcommand{\bits}{\{-1,1\}}
\newcommand{\zo}{\{0,1\}}
\newcommand{\bn}{\bits^n}
\newcommand{\isafunc}{:\bn\to\bits}
\newcommand{\isazofunc}{:\zo^n\to\zo}
\newcommand{\sumi}{\sum_{i=1}^n}

\newcommand{\Fsm}{\mathbb{F}_{\mathrm{sm}}}
\newcommand{\Ksm}{\mathbb{K}_{\mathrm{sm}}}
\newcommand{\lfrac}{\textstyle \frac}



% Overline
\newcommand{\ol}[1]{\overline{#1}}

% BOLD
\newcommand{\ind}{{\bf 1}}

\newcommand{\bA}{\mathbf{A}}
\newcommand{\bB}{\mathbf{B}}
\newcommand{\bC}{\mathbf{C}}
\newcommand{\bD}{\mathbf{D}}
\newcommand{\bE}{\mathbf{E}}
\newcommand{\bF}{\mathbf{F}}
\newcommand{\bG}{\mathbf{G}}
\newcommand{\bH}{\mathbf{H}}
\newcommand{\bI}{\mathbf{I}}
\newcommand{\bJ}{\mathbf{J}}
\newcommand{\bK}{\mathbf{K}}
\newcommand{\bL}{\mathbf{L}}
\newcommand{\bM}{\mathbf{M}}
\newcommand{\bN}{\mathbf{N}}
\newcommand{\bO}{\mathbf{O}}
\newcommand{\bP}{\mathbf{P}}
\newcommand{\bQ}{\mathbf{Q}}
\newcommand{\bR}{\mathbf{R}}
\newcommand{\bS}{\mathbf{S}}
\newcommand{\bT}{\mathbf{T}}
\newcommand{\bU}{\mathbf{U}}
\newcommand{\bV}{\mathbf{V}}
\newcommand{\bW}{\mathbf{W}}
\newcommand{\bX}{\mathbf{X}}
\newcommand{\bY}{\mathbf{Y}}
\newcommand{\bZ}{\mathbf{Z}}


\newcommand{\boldf}{\boldsymbol{f}} % \bf taken
\newcommand{\ba}{\boldsymbol{a}}
\newcommand{\bb}{\boldsymbol{b}}
\newcommand{\be}{\boldsymbol{e}}
\newcommand{\bg}{\boldsymbol{g}}
\newcommand{\bh}{\boldsymbol{h}}
\newcommand{\bj}{\boldsymbol{j}}
\newcommand{\bk}{\boldsymbol{k}}
\newcommand{\bm}{\boldsymbol{m}}
\newcommand{\bp}{\boldsymbol{p}}
\newcommand{\boldr}{\boldsymbol{r}} % \br taken
\newcommand{\bs}{\boldsymbol{s}}
\newcommand{\bt}{\boldsymbol{t}}
\newcommand{\bu}{\boldsymbol{u}}
\newcommand{\bv}{\boldsymbol{v}}
\newcommand{\bw}{\boldsymbol{w}}
\newcommand{\bx}{\boldsymbol{x}}
\newcommand{\by}{\boldsymbol{y}}
\newcommand{\bz}{\boldsymbol{z}}

\newcommand{\boldR}{\boldsymbol{R}}


\newcommand{\bsigma}{\boldsymbol{\sigma}}
\newcommand{\bpi}{\boldsymbol{\pi}}
\newcommand{\bell}{\boldsymbol{\ell}}
\newcommand{\bnu}{\boldsymbol{\nu}}
\newcommand{\bxi}{\boldsymbol{\xi}}
\newcommand{\btau}{\boldsymbol{\tau}} 
\newcommand{\brho}{{\boldsymbol{\rho}} }
\newcommand{\bgamma}{{\boldsymbol{\gamma}} }
\newcommand{\bphi}{{\boldsymbol{\varphi}} }


\newcommand{\bvsigma}{{\boldsymbol{\varsigma}} }
\newcommand{\bvrho}{{\boldsymbol{\varrho}}}

\newcommand{\bGamma}{{\mathbf{\Gamma}} }


% Calligraphy

\newcommand{\calA}{\mathcal{A}}
\newcommand{\calB}{\mathcal{B}}
\newcommand{\calC}{\mathcal{C}}
\newcommand{\calD}{\mathcal{D}}
\newcommand{\calE}{\mathcal{E}}
\newcommand{\calF}{\mathcal{F}}
\newcommand{\calG}{\mathcal{G}}
\newcommand{\calH}{\mathcal{H}}
\newcommand{\calI}{\mathcal{I}}
\newcommand{\calJ}{\mathcal{J}}
\newcommand{\calK}{\mathcal{K}}
\newcommand{\calL}{\mathcal{L}}
\newcommand{\calM}{\mathcal{M}}
\newcommand{\calN}{\mathcal{N}}
\newcommand{\calO}{\mathcal{O}}
\newcommand{\calP}{\mathcal{P}}
\newcommand{\calQ}{\mathcal{Q}}
\newcommand{\calR}{\mathcal{R}}
\newcommand{\calS}{\mathcal{S}}
\newcommand{\calT}{\mathcal{T}}
\newcommand{\calU}{\mathcal{U}}
\newcommand{\calV}{\mathcal{V}}
\newcommand{\calW}{\mathcal{W}}
\newcommand{\calX}{\mathcal{X}}
\newcommand{\calY}{\mathcal{Y}}
\newcommand{\calZ}{\mathcal{Z}}

\newenvironment{solution}
  {\begin{proof}[\textbf{Solution}]}
  {\end{proof}}
 
\usepackage{cleveref}
\crefname{originalsection}{original section}{original sections}
\Crefname{originalsection}{Original section}{Original sections}
\makeatletter
\newlabel{subsec:our-results}{{1.3}{source}{}{sourcepointer-1}{}}
\newlabel{subsec:our-results@cref}{{[originalsection][0][]1.3}{[1][0][]source}{}{}{}}
\makeatother
\title{Near-optimal separation of ordered from unordered set-multilinear algebraic branching programs}
\author{Prerona Chatterjee, Deepanshu Kush, Shubhangi Saraf, Amir Shpilka}
\date{doi:10.4230/LIPIcs.CCC.2024.20; CC BY-NC-SA4.0}
\begin{document}
\maketitle
\noindent\textit{Editorial note.} Original source: Prerona Chatterjee, Deepanshu Kush, Shubhangi Saraf and Amir Shpilka; arXiv:2312.15874v3 (2024-02-19), https://arxiv.org/abs/2312.15874v3. License: Creative Commons Attribution-NonCommercial-ShareAlike4.0, https://creativecommons.org/licenses/by-nc-sa/4.0/.\par\medskip
\noindent\textit{Editorial note.} Non-commercial research collection. Changes: complete author theorem/proof/context excerpts and original macros/bibliography in standalone layout; unrelated sections and frontmatter omitted. No mathematical corrections or additional proofs. Adapted presentation retains CC BY-NC-SA4.0 including NC/SA restrictions.\par\medskip
\noindent\textit{Allowed standard prerequisites (editorial).}\begin{itemize}
\item Nisan ordered-ABP rank-width lower bound. Author Lemma lem:nisan-abpsize explicitly cites Nisan91 and states the exact relative-rank form inside the full proof.
\item Standard relative-rank multiplicativity and inner-product rank gadget. Author Claim clm:rk-props and cited Observation obs:ip-full-rank; exact statements retained.
\end{itemize}\medskip
\ifdefined\hypertarget\hypertarget{sourcepointer-1}{}\fi
\noindent\textit{Source reference.} Original section 1.3: number-theory/sources/ccc-2024-20/native/main.tex, source lines 150--150. Optional introductory pointer to results; counted full target statements and proofs are included.\par\medskip
An \emph{algebraic circuit} over a field $\F$ for a multivariate polynomial $P(x_1,\ldots,x_N)$ is a directed acyclic graph (DAG)
whose internal vertices (called gates) are labeled as either $+$ (sum) or $\times$ (product), and leaves
(vertices of in-degree zero) are labeled by the variables $x_i$ or constants from $\F$. 
A special output gate (the root of the DAG) represents the polynomial
$P$. 
If the DAG happens to be a tree, such a resulting circuit is called an \emph{algebraic formula}. 
The size of a circuit or formula is the number of nodes in the DAG. We also consider the product-depth of the circuit, which is the maximum number of product gates on a root-to-leaf path. 
The class $\VP$ (respectively, $\VF$) is then defined to be the collection of all polynomials having at most polynomially large degree which can be computed by polynomial-sized circuits (respectively, formulas).

An \emph{algebraic branching program} (ABP) is a layered DAG with two special nodes in
it: a start-node and an end-node. All edges of the ABP go from layer $\ell -1$ to layer $\ell$ for some $\ell$ (say start-node is the unique node in layer $0$ and end-node is the unique node in the last layer) and are labeled by a linear polynomial. Every directed path $\gamma$ from start-node to end-node computes the monomial $P_\gamma$,
which is the product of all labels on the path $\gamma$. The ABP computes the polynomial $P = \sum_{\gamma} P_\gamma$,
where the sum is over all paths $\gamma$ from start-node to end-node. Its size is simply the number of nodes in the DAG, its \emph{depth} is the length of the longest path from the start-node to the end-node, and \emph{width} is the maximum number of nodes in any layer. The class $\VBP$ is then defined to be the collection of all polynomials (with polynomially-bounded degree) which can be computed by polynomial-sized branching programs. ABPs are known to be of intermediate complexity between formulas and circuits; in other words, we know the inclusions $\VF\subseteq\VBP \subseteq \VP\sse \VNP$.

A polynomial is said to be \emph{homogeneous}
if each monomial has the same total degree and \emph{multilinear} if every variable occurs
at most once in any monomial. Now, suppose that the underlying variable set is partitioned into $d$ sets
$X_1,\ldots, X_d$. Then the polynomial is said to be \emph{set-multilinear} with respect to this variable partition if each
monomial in $P$ has \emph{exactly} one variable from each set. Note that a set-multilinear polynomial is both multilinear and homogeneous, and has degree precisely $d$ if it is set-multilinear over $d$ sets. Next, we define different models
of computation corresponding to these variants of polynomials classes. An algebraic formula/branching program/circuit is set-multilinear with respect to a variable partition
$(X_1,\ldots, X_d)$ if each internal node in the formula/branching program/circuit computes a set-multilinear polynomial.\footnote{Of course, a non-root node need not be set-multilinear with respect to the \emph{entire} variable partition. Nevertheless, here we demand that it must be set-multilinear with respect to some \emph{subset} of the collection $\{X_1,\ldots,X_d\}$.} Multilinear and homogeneous formulas/branching programs/circuits are defined analogously.  

\begin{definition}[Ordered smABP]\label{def:smABP}
Given a variable partition $(X_1,\ldots, X_d)$, we say that a set-multilinear branching program of depth $d$ is said to be \emph{ordered} with respect to an 
ordering (or permutation) $\sigma\in S_d$ if for each $\ell\in[d]$, all edges of the ABP from layer $\ell-1$ to layer $\ell$ are labeled using a linear form over the variables in $X_{\sigma(\ell)}$. It is simply said to be \emph{ordered} if there \emph{exists} an ordering $\sigma$ such that it is ordered with respect to $\sigma$.
\end{definition}

At this point, it is essential to take note of the terminology in this context: in this paper, a general (or ``unordered'') set-multilinear branching program refers to an ABP for which each internal node computes a polynomial that is set-multilinear with respect to \emph{some} subset of the global partition, whereas an \emph{ordered} set-multilinear branching program is more specialized and has the property that any two nodes in the same layer compute polynomials that are set-multilinear with respect to the \emph{same} partition.

\begin{definition}[$\sbp$]\label{def:sum}
Given a polynomial $P(X)$ that is set-multilinear with respect to the variable partition $X = (X_1,\ldots, X_d)$, we say that $\sum_{i=1}^t A_i$ is a $\sbp$ computing $P$ if indeed $\sum_{i=1}^t A_i(X) = P(X)$, and
each $A_i$ is an ordered set-multilinear branching program i.e., each $A_i$ is ordered with respect to some $\sigma_i\in S_d$. We call $t$ (i.e., the number of summands in a $\sbp$) its \emph{support size} and define its \emph{max-width} and \emph{total-width} to be the maximum over the width of each $A_i$ and the sum of the width of each $A_i$, respectively.    
\end{definition}

\begin{theorem}[Near-Optimal Separation between Ordered and Unordered smABPs]\label{thm:sep-ord-unord}
    There is a polynomial $G_{n,d}(X)$ which is set-multilinear over the variable partition $X = (X_1,\ldots,X_d)$ with
each $|X_i| \leq n$, and such that:
\begin{itemize}
    \item it has a set-multilinear branching program of size $\poly(n,d)$,
    \item but any ordered set-multilinear branching program computing $G_{n,d}$ requires width $n^{\Omega(d)}$. 
\end{itemize}
\end{theorem}

\section{Preliminaries}\label{sec:prel}
\subsection{Relative Rank and its Properties}

We first describe the notation that we need to define the measures that we use to prove our results described in Section \ref{subsec:our-results}. Instead of directly working with the rank of the partial derivative matrix, we work with the following normalized form.

\begin{definition}\label{def:rk-rest}
Let $w = (w_1, w_2,\ldots, w_d)$ be a tuple (or word) of non-zero real numbers. For a subset $S \subseteq [t]$, we shall refer to the sum $\sum_{i\in S} w_i$ by $w_S$,  and by $w|_S$, we will refer to the tuple obtained by considering only the
elements of $w$ that are indexed by $S$. 
Given a word $w=(w_1,\ldots,w_d)$, we denote by $\ol{X}(w)$ a tuple of $d$ sets of variables $(X(w_1),\ldots,X(w_d))$ where
$|X(w_i)| = 2^{|w_i|}$.\footnote{In particular, $2^{|w_i|}\in\mathbb{N}$.}
We denote by $\Fsm[\calT]$ the set of set-multilinear polynomials over the tuple
of sets of variables $\calT$.
\end{definition}


\begin{definition}[Relative Rank Measure of \cite{LST1}]\label{def:relrk}
Let $\ol{X}=(X_1,\ldots,X_d)$ be a tuple of sets of variables such that $|X_i|=n_i$ and let $f\in \Fsm[\ol{X}]$.
 Let $w = (w_1, w_2,\ldots, w_d)$ be a tuple (or word) of non-zero real numbers such that $2^{|w_i|}  = n_i $ for all $i \in [d]$. Corresponding to a word $w$, define $\calP_w \coloneqq \{i\ |\ w_i > 0\}$ and $\calN_w \coloneqq \{i\ |\ w_i < 0\}$. Let $\calM^{\calP}_{w}$ be the
set of all set-multilinear monomials over the subset of the variable sets $X_1, X_2,\ldots, X_d$ precisely
indexed by $\calP_w$, and similarly let $\calM^{\calN}_{w}$ be the set of all set-multilinear monomials over these variable
sets indexed by $\calN_w$.

Define the ‘partial derivative matrix’ matrix $\calM_w(f)$ whose rows are indexed by the elements
of $\calM^{\calP}_w$ and columns indexed by the elements of $\calM^{\calN}_{w}$ as follows: the entry of this matrix corresponding to a
row $m_1$ and a column $m_2$ is the coefficient of the monomial $m_1\cdot m_2$ in $f$. We define
\[
\rk_w(f) \coloneqq \frac{\mathrm{rank}(\calM_w(f))}{\sqrt{|\calM^{\calP}_w|\cdot |\calM^{\calN}_w|}} = \frac{\mathrm{rank}(\calM_w(f))}{2^{\frac{1}{2}\sum_{i\in[d]}|w_i|}}.
\]
\end{definition}


The following is a simple result that establishes various useful properties of the relative rank measure.
\begin{claim}[\cite{LST1}]\label{clm:rk-props}
\begin{enumerate}
    \item(Imbalance) Say $f \in \Fsm[\ol{X}(w)]$. Then, $\rk_w(f)\leq 2^{-|w_{[d]}|/2}$.
    \item(Sub-additivity) If $f,g \in \Fsm[\ol{X}(w)]$, then $\rk_w(f+g)\leq \rk_w(f)+\rk_w(g)$.
    \item(Multiplicativity) Say $f = f_1 f_2\cdots f_t $ and assume that for each $i\in [t]$, $f_i\in \Fsm[\ol{X}(w|_{S_i})]$, where $(S_1, \ldots, S_t)$ is a partition of $[d]$. Then
    \[
    \rk_w(f) = \prod_{i\in[t]} \rk_{w|_{S_i}}(f_i).
    \]
\end{enumerate}
\end{claim}

We will repeatedly make use of the following.

\begin{theorem}[Chernoff Bound, as stated in \cite{mitzenmacher_upfal_2005}]\label{thm:Chernoff}
Suppose $X_1, \ldots, X_n$ are independent random variables taking values in $\zo$. Let $X$ denote their sum and let $\mu = \E[X]$ denote the expected value of the sum. Then for any $\delta > 0$,
\begin{eqnarray*} \Pr[X\ge (1+\delta)\mu] &\le& \exp\left(-\frac{\delta^2\mu}{2+\delta}\right), \text{ if }0\leq \delta, \\
  \Pr[X\le (1-\delta)\mu] &\le& \exp\left(-\frac{\delta^2\mu}{2}\right), \text{ if }0< \delta <1, \\ 
  \Pr[|X-\mu|\geq \delta \mu]&\leq& 2\cdot \exp\left(-\frac{\delta^2\mu}{3}\right),\text{ if }0< \delta <1.\\
\end{eqnarray*}
\end{theorem}

\subsection{Inner Product Gadget}
The following observation is used crucially to construct the hard polynomials in $\VP$ as well as $\VBP$. 
\begin{obs}[\cite{KushS23}]\label{obs:ip-full-rank}
    Let $n = 2^k$ and $X_1 = \{x_{1,1},\ldots,x_{1,n}\}$ and $X_2 = \{x_{2,1},\ldots,x_{2,n}\}$ be two disjoint sets of variables. Then, for any \emph{symmetric} word $w\in\{k,-k\}^2$ (i.e., where $w_1 + w_2 = 0$) and for the inner product `gadget' $ f = X_1\cdot X_2 = \sumi x_{1,i}x_{2,i}$, $\rk_w(f) = 1$ i.e., $\calM_w(f)$ is full-rank.
\end{obs}


\subsection{A Hard Set-Multilinear Polynomial in $\VBP$}\label{subsec:VBP}
\subsubsection{Arc-partition Measure Description}
This subsection is adapted from Section 2 of \cite{DvirMPY12}. Let $n = 2^k$, $d\leq n$ be an even integer, and let $X = (X_1, X_2, \ldots, X_{d})$ be a collection of disjoint sets of $n$ variables each.
An \emph{arc-partition} will be a special kind of \emph{symmetric} word $w\in \{-k,k\}^d$ (i.e., a one-to-one map $\Pi$ from $X$ to $\{-k,k\}^d$). For the purpose of this subsection, the reader can even choose to think of the alphabet of $w$ as $\{-1,1\}$ (i.e., one `positive' and one `negative' value) -- we use $k,-k$ only to remain consistent with \autoref{def:relrk}.

Identify $X$ with the set $\{1, 2,\ldots, d\}$ in the natural way. Consider the $d$-cycle
graph, i.e., the graph with nodes $\{1, 2,\ldots, d\}$ and edges between $i$ and $i + 1$ modulo $d$. For
two nodes $i \neq j$ in the $d$-cycle, denote by $[i, j]$ the arc between $i, j$, that is, the set of nodes on the
path $\{i, i + 1, \ldots , j - 1, j\}$ from $i$ to $j$ in $d$-cycle.
First, define a distribution $\calD_P$ on a family of pairings (a list of disjoint pairs of nodes in the
cycle) as follows. A random pairing is constructed in $d/2$ steps. At the end of step $t \in [d/2]$, we
shall have a pairing $(P_1,\ldots , P_t)$ of the arc $[L_t, R_t]$. The size of $[L_t
, R_t]$ is always $2t$. The first pairing
contains only $P_1 = \{L_1, R_1\}$ with $L_1 = 1$ and $R_1 = 2$. Given $(P_1, \ldots, P_t)$ and $[L_t, R_t]$, define the
random pair $P_t+1$ (independently of previous choices) by 

$$
P_{t+1} = \begin{cases}
    \{L_t - 2, L_t - 1\} &\text{ with probability $1/3$}\\
    \{L_t - 1, R_t + 1\} &\text{ with probability $1/3$}\\
    \{R_t + 1, R_t + 2\} &\text{ with probability $1/3$}\\
\end{cases}
$$

Define
$$
[L_{t+1},R_{t+1}] = [L_t,R_t] \cup P_{t+1}.
$$

So, $L_{t+1}$ is either $L_t - 2$, $L_t - 1$ or $L_t$, each value is obtained with probability $1/3$, and similarly
(but not independently) for $R_{t+1}$.

The final pairing is
$P = (P_1, P_2, \ldots , P_{d/2}
)$.
Denote by $P \sim \calD_P$ a pairing distributed according to $\calD_P$.

Once a pairing $P$ has been obtained, a word $w\in \{-k,k\}^d$ is obtained by simply randomly assigning $+k$ and $-k$ to the indices of any pair $P_i$. More formally, for every
$t \in [d/2]$, if $P_t = \{i_t
, j_t\}$, let with probability $1/2$, independently of all other choices,
$$
w_{i_t} = +k \text{ and } w_{j_t} = -k,
$$
and with probability 1/2,
$$
w_{i_t} = -k \text{ and } w_{j_t} = +k.
$$
Denote by $w\sim \calD$ a word in $\{-1,1\}^n$ that is sampled using this procedure. We call such a word an \emph{arc-partition}. For a pair $P_t = \{i_t,j_t\}$, we refer to $i_t$ and $j_t$ as \emph{partners}. 

\begin{definition}[Arc-full-rank]\label{def:arc-full-rk}
We say that a polynomial $f$ that is set-multilinear over $X = (X_1, \ldots, X_{d})$ is \textbf{arc-full-rank} if for every arc-partition $w\in \{-k,k\}^d$, $\rk_w(f) = 1$.
\end{definition}


\subsubsection{Construction of an Arc-full-rank Polynomial}\label{subsubsec:arc-full-rank-poly}

Below, we describe a simple construction of a polynomial sized ABP that computes an arc-full-rank set-multilinear polynomial.
The high-level idea is to construct an ABP in which every path between start-node and end-node
corresponds to a specific execution of the random process which samples arc-partitions. Each node in the ABP corresponds to an arc $[L, R]$, which sends an edge to each of the nodes $[L - 2, R], [L -
1, R + 1]$ and $[L, R + 2]$. The edges have specially chosen labels that help guarantee full rank with respect to
every arc-partition. For simplicity of presentation, we allow the edges of the program to be labeled
by degree four set-multilinear polynomial polynomials over the corresponding subset of the variable partition. This assumption can be easily removed by replacing
each edge with a polynomial-sized ABP computing the corresponding degree four polynomial.

Formally, the nodes of the program are even-size arcs in the $d$-cycle, $d$ an even integer. The
start-node of the program is the empty arc $\emptyset$ and the end-node is the whole cycle $[d]$ (both are
“special” arcs). Let $X = (X_1,\ldots, X_{d})$ be a collection of sets of variables where each $|X_i| = n$, and similarly, let $Y = (Y_1,\ldots, Y_{d})$ be a distinct collection of sets of variables where each $|Y_i| = n$ (we shall refer to the $Y$-variables as \emph{auxiliary} variables). For $i$ and $j$ in $\{1,\ldots,d\}$, let $X_i \cdot X_j$ denote the inner-product quadratic form $\sum_{k = 1}^n x_{ik}x_{jk}$. Here, we shall assume that $X_i = \{x_{i,1},\ldots, x_{i,n}\}$ and $Y_i = \{y_{i,1},\ldots, y_{i,n}\}$.

Construct the branching program by connecting a node/arc of size $2t$ to three nodes/arcs of
size $2t + 2$. For $t = 1$, there is just one node $[1, 2]$, and the edge from start-node to it is labeled
$y_{1,2}y_{2,1}(X_0\cdot X_1)$. For $t > 1$, the node $[L, R] \supset [1, 2]$ of size $2t < d$ is connected to the three nodes:
$[L-2, R] , [L-1, R+ 1]$, and $[L, R+ 2]$. (It may be the case that the three nodes are the end-node.)
The edge labeling is: 
\begin{itemize}
    \item The edge  between $[L, R]$ and $[L-2, R]$ is labeled $y_{L-2,L-1}y_{L-1,L-2}(X_{L-2}\cdot X_{L-1})$.
    \item The
edge between $[L, R]$ and $[L - 1, R + 1]$ is labeled $y_{L-1,R+1}y_{R+1,L-1}(X_{L-1}\cdot X_{R+1})$.
    \item The edge between $[L, R]$
and $[L, R + 2]$ is labeled $y_{R+1,R+2}y_{R+2,R+1}(X_{R+1}\cdot X_{R+2})$.
\end{itemize}
  
Consider the ABP thus described, and the polynomial $G_{n,d}$ it computes. For every path $\gamma$
from start-node to end-node in the ABP, the list of edges along $\gamma$ yields a pairing $P$; every edge $e$
in $\gamma$ corresponds to a pair $P_e = \{i_e, j_e\}$ of nodes in $d$-cycle. Thus,

\begin{equation}\label{eq:poly-def}
  G_{n,d} = \sum_{\gamma} \prod_{e = \{i_e, j_e\} \in \gamma}y_{i_e,j_e}y_{j_e,i_e} \cdot (X_{i_e}\cdot X_{j_e}). 
\end{equation}

where the sum is over all paths $\gamma$ from start-node to end-node. 

\begin{remark}\label{rem:parings-paths}
There is in fact a one-to-one
correspondence between pairings $P$ and such paths $\gamma$ (this follows by induction on $t$). {Note that this is true only because pairings are tuples i.e., they are \emph{ordered} by definition. Otherwise, it is of course still possible to obtain the same \emph{set} of pairs in a given pairing using multiple different orderings.} The sum
defining $G_{n,d}$ can be thought of, therefore, as over pairings $P$.  
\end{remark}

The following statement summarizes the main useful property of $G_{n,d}$.

\begin{lem}[\cite{KushS23}]\label{lem:Fn-VBP-props}
Over any field $\F$ of characteristic zero, the polynomial $ G_{n,d}$ defined above is arc-full-rank as a set-multilinear polynomial in the variables $X$ over the field $\F(Y)$ of rational functions
in $Y$.
\end{lem}

\begin{proof}
    Let $w\sim \calD$ be an arc-partition. 
    We want to show that $\calM_w(G_{n,d})$ has full rank. 
    The arc-partition $w$ is defined from a pairing $P = (P_1,\ldots,P_{d/2})$ (though as discussed in \autoref{rem:parings-paths}, there could be multiple such $P$). 
    The pairing $P$ corresponds to a path $\gamma$ from start-node to end-node. 
    Consider the polynomial $f$ that is obtained by setting every $y_{i,j} = y_{j,i} = 0$ in $F$ such that $\{i,j\}$ is not a pair in $P$, and setting every $y_{i,j} = y_{j,i} = 1$ for every pair $\{i,j\}$ in $P$. 
    Then, it is easy to see that the only terms that survive in \autoref{eq:poly-def} correspond to paths (and in turn, pairings) which have the same underlying \emph{set} of pairs as $P$.
    As a consequence, $f$ is simply some non-zero constant times a polynomial which is full-rank (recall \autoref{obs:ip-full-rank}).
    $M_w(f)$ being full rank then implies that $M_w(G_{n,d})$ is also full-rank.
\end{proof}

\section{Optimal Separation between Ordered and Unordered smABPs}\label{sec:opt-sep-sml-osml}
In this section, we prove \autoref{thm:sep-ord-unord}. The result is that $G_{n,d}(X_1,\ldots, X_d,Y_1,\ldots,Y_d)$, as described in Section \ref{subsec:VBP}, does not have small \emph{ordered} set-multilinear branching programs with respect to \emph{any} ordering of the $X_i$s and $Y_j$s. 
More precisely, we claim that any set-multilinear branching program computing $G_{n,d}$ which is ordered with respect to some permutation of the sets in the collection $\calS = \{X_1,\ldots, X_d,Y_1,\ldots,Y_d\}$ must have size at least $n^{\Omega(d)}$.

\begin{proof}[Proof of \autoref{thm:sep-ord-unord}]
    Let $A$ be a set-multilinear branching program computing $G_{n,d}$, which is ordered with respect to a permutation $\pi$ of the sets in $\calS$.
    That is $\pi: [2d]\rightarrow \calS$ is a bijective map where $\pi(i) = Z$, for some $Z\in \calS$, if all edges between layer $i-1$ and layer $i$ are labelled using linear forms in the variables of $Z$. 
    Let $\pi_X:[d]\rightarrow \{X_1,\ldots,X_d\}$ be the ordering inherited from $\pi$ by the $X$ sets, that is $\pi_X(i) = X_j$ if $X_j$ is the $i$-th $X$-set appearing in the sequence $\pi(1),\ldots,\pi(2d)$. 
    Define a word $w\in \{-k,k\}^d$ using $\pi_X$ as follows: let 
    \[
        w_i = k \text{ if and only if } \pi_X^{-1}(X_i)\in [d/2].
    \] 
    That is, the edges labelled using linear forms in $X_i$ appear on the `left' half in $A$, when the $Y$ sets are ignored. 
    Note that $w$ is \emph{symmetric}, or in other words, $w_{[d]} = 0$.

    Now, for $i\in[d]$, define $s_i(A)$ to be the size of layer $\pi^{-1}(\pi_X(i))$ in $A$.
    The technique used by Nisan \cite{Nisan91} can be used directly to show the following.

    \begin{lem}[\cite{Nisan91}]\label{lem:nisan-abpsize}
        If $G_{n,d}(X_1,\ldots,X_d)\in \mathbb{K}[X_1,\ldots,X_d]$ is thought of as a polynomial only over the $X$ variables (here $\mathbb{K} = \F(Y)$, the field of rational functions in $Y$), then $n^{d/2}\cdot \rk_w(G_{n,d})\leq s_{d/2}(A)$.\footnote{The $n^{d/2}$ factor comes from using relative-rank rather than rank.}
    \end{lem}

    Thus, in order to prove a $n^{\Omega(d)}$ lower bound on the width of $A$, it suffices to lower bound $\rk_w(G)$ by $n^{-\alpha d}$ for some $0 < \alpha < 1/2$.
    We prove such a lower bound using the following lemma.

    \begin{lem}\label{lem:overlap}
        Given any $w\in\{-k,k\}^d$ with $w_{[d]} = 0$, there exists an arc-partition $v\in\{-k,k\}^d$ obtained from a pairing $P = (P_1,\ldots,P_{d/2})$ such that $w$ `splits' a constant fraction of the pairs in $P$: more precisely, there is a set $S\subset P$ of size at least $d/8$ such that if $(i,j) \in S$, then $w_i + w_j = 0$.
    \end{lem}

    Let us first prove the lower bound on $\rk_w(G_{n,d})$ assuming this lemma. 
    %Given the pairing $P = (P_1,\ldots,P_{d/2})$ from \autoref{lem:overlap}, note that it uniquely corresponds to a path $\gamma$ from start-node to end-node. 
    Consider the polynomial $f$ that is obtained by setting every $y_{i,j} = y_{j,i} = 0$ in $G_{n,d}$ such that $\{i,j\}$ is not a pair in $P$, and setting every $y_{i,j} = y_{j,i} = 1$ for every pair $\{i,j\}$ in $P$ (where $P$ is as in \autoref{lem:overlap}). 
    Then, it is easy to see that the only terms that survive in \autoref{eq:poly-def} correspond to paths (and in turn, pairings) which have the same underlying \emph{set} of pairs as $P$. 
    As a consequence,
    \[
        f(X)= c_P \cdot \prod_{t = 1}^{d/2} (X_{i_t}\cdot X_{j_t})
    \]
    for a non-zero constant $c_P$.
    Next, notice that the rank of the matrix $M_w(f)$ (over the field $\F$) serves as a lower bound on the rank of the matrix $M_w(G_{n,d})$ (over the field $\F(Y)$).
    Indeed, if the former is $r$ then there is a non-vanishing $r\times r$ minor of $M_w(f)$. 
    But this implies that the determinant of the corresponding $r\times r$ sub-matrix in $M_w(G_{n,d})$ must be a non-zero polynomial in $\F[Y]$ as it has a non-zero evaluation.
    From this we conclude that $M_w(G_{n,d})$ also has a non-vanishing $r\times r$ minor and therefore has rank at least $r$. 
    Stated in terms of relative rank, we have showed that $\rk_w(f)\leq \rk_w(G_{n,d})$.

    Now, given the set $S\subset [d/2]$ from \autoref{lem:overlap}, we can lower bound the relative rank of $M_w(f)$ as follows: first, note that if for some $t\in[d/2]$, the pair $P_t$ of $P$ is in $S$, then by \autoref{obs:ip-full-rank}, $\rk_{w|_{P_t}}(X_{i_t}\cdot X_{j_t}) = 1$. 
    Secondly, if $w$ does \emph{not} split some pair $P_t$ of $P$ (i.e., $w_{i_t} = w_{j_t}$), then $M_{w|_{P_t}}(X_{i_t}\cdot X_{j_t})$ is simply a one-dimensional vector (either $1\times n^2$ or $n^2\times 1$) and we trivially have $\rk_{w|_{P_t}}(X_{i_t}\cdot X_{j_t}) = 1/n$. 
    Combining, and using the multiplicativity of $\rk_w(\cdot)$ (third item of \autoref{clm:rk-props}), 
    \begin{align*}
        \rk_w(f(X)) = \rk_w(f(X)/c_P) &= \prod_{t = 1}^{d/2} \rk_{w|_{P_t}}(X_{i_t}\cdot X_{j_t})\\
        &= \prod_{t\in S} \underbrace{\rk_{w|_{P_t}}(X_{i_t}\cdot X_{j_t})}_{= 1} \prod_{t\in[d/2]\setminus S}\underbrace{\rk_{w|_{P_t}}(X_{i_t}\cdot X_{j_t})}_{\geq 1/n}
        \geq \frac{1}{n^{d/2 - |S|}}.
    \end{align*}
    
    Therefore, we have $\rk_w(G_{n,d})\geq \rk_w(f) \geq 1/n^{d/2 - |S|}$ and hence, by \autoref{lem:nisan-abpsize}, the size of $A$ is at least $s_{d/2}(A)\geq n^{d/2}\cdot \rk_w(G_{n,d}) \geq n^{|S|} = n^{\Omega(d)}$.
\end{proof}

We now give the proof of \autoref{lem:overlap}.

\begin{proof}[Proof of \autoref{lem:overlap}]
    Let $L_+, L_-, R_+, R_- \in [d]$ be defined as follows.
    \begin{align*}
        R_+ &= \left\{i \in  \left[ 3, \frac{d}{2}+1 \right] \quad : \quad w_i > 0 \right\}
        \qquad \qquad \qquad
        R_- = \left\{i \in  \left[ 3, \frac{d}{2}+1 \right] \quad : \quad w_i < 0 \right\}\\
        L_+ &= \left\{i \in  \left[\frac{d}{2}+2, d \right] \quad : \quad w_i > 0 \right\}
        \qquad \qquad \qquad
        L_- = \left\{i \in  \left[\frac{d}{2}+2, d \right] \quad : \quad w_i < 0 \right\}
    \end{align*}
    Clearly, either $|R_+|\geq\frac{d}{4}$ or $|R_-|\geq\frac{d}{4}$.
    Without loss of generality, let us assume that $|R_+|\geq\frac{d}{4}$.
    
    Also, without loss of generality\footnote{If this is not the case, then $|L_-| -2 \leq |R_+| \leq |L_-| + 2$ and we redefine $R_+, L_-$ to be the largest possible subsets of the originally defined $R_+, L_-$, respectively, such that $|R_+| = |L_-|$.} we can assume that $w_1 + w_2 = 0$.
    Then $|R_+| = |L_-|$, and say $|R_+| = p$.
    Further, let 
    \[
        R_+ = \{i_1, \ldots, i_p\} \text{ with } i_1 < \cdots < i_p \quad \text{ and } \quad L_- = \{j_1, \ldots, j_p\} \text{ with } j_1 > \cdots > j_p.
    \]
    We can then define an initial pairing, $P^0 = \{(i_\ell, j_\ell)\}_{\ell \in [p]}$.
    Let us also define the set $S^0 = \emptyset$.
    The goal is to iteratively update $P^0$ such that, at the end, we have a pairing corresponding to an arc-partition. 
    We will also update $S^0$ at each step so that, at the end, each pair in $S$ are of opposite signs.
    
    Let $u_0=2, v_0=1$.
    Intuitively, $u_\ell, v_\ell$ are the right most and left-most points, respectively, of the partial arc-partition, $P^\ell$, defined till the $\ell$-th iteration.
    Given $u_{\ell-1}, v_{\ell-1}, P^{\ell-1}, S^{\ell -1}$ for any $\ell \in [p]$, we will define $u_\ell, v_\ell, P^\ell, S^\ell$ as follows.
    Note that the calculations are $\pmod d$ when $\ell=1$.
    \begin{enumerate}
        \item[\underline{\textbf{Case 1:}} ]  $(i_\ell - u_{\ell-1})$ and $(v_{\ell-1} - j_\ell)$ are both odd.
        \[
            P^\ell = P^{\ell-1} \cup \{(u_{\ell-1}+1, u_{\ell-1}+2), \ldots, (i_\ell -2, i_\ell -1)\} \cup \{(v_{\ell-1}-2, v_{\ell-1}-1), \ldots, (j_\ell +1, j_\ell +2)\}
        \]
        \[
            u_\ell = i_\ell \qquad \qquad v_\ell = j_\ell \qquad \text{ and } \qquad S^\ell = S^{\ell -1} \cup \{(i_\ell, j_\ell) \}
        \]
        \item[\underline{\textbf{Case 2:}} ] $(i_\ell - u_{\ell-1})$ is even and $(v_{\ell-1} - j_\ell)$ is odd.
        
        We first define
        \[
            Q^\ell = \begin{cases}
                (P^{\ell-1} \setminus \{(i_\ell, j_\ell)\}) & \cup \quad \{(u_{\ell-1}+1, u_{\ell-1}+2), \ldots, (i_\ell -1, i_\ell)\} \\
                & \cup \quad \{(v_{\ell-1}-2, v_{\ell-1}-1), \ldots, (j_\ell +1, j_\ell +2\}, 
                \qquad \text{if } (i_\ell - u_{\ell-1}) > 0\\
                (P^{\ell-1} \setminus \{(i_\ell, j_\ell)\}) &\cup \quad \{(v_{\ell-1}-2, v_{\ell-1}-1), \ldots, (j_\ell +1, j_\ell +2)\} 
                \qquad \text{ if } (i_\ell - u_{\ell-1}) = 0,
            \end{cases}
        \]
        
        and then define
        \[
            P^\ell = \begin{cases}
                Q^\ell \cup \{(j_\ell -1, j_\ell)\} \quad \text{ if } \quad j_\ell -1 \in L_+.\\
                Q^\ell \cup \{(i_\ell +1, j_\ell)\} \quad \text{ otherwise.} 
                \end{cases} 
        \]
        
        Also,
        \[
            u_\ell = \begin{cases}
                i_\ell \qquad \quad \text{ if } \quad j_\ell -1 \in L_+\\
                i_\ell +1 \quad \text{ otherwise.} 
            \end{cases} 
            \qquad \qquad 
            v_\ell = \begin{cases}
                j_\ell - 1 \quad \text{ if } \quad j_\ell -1 \in L_+\\
                j_\ell \qquad \quad \text{otherwise.} 
            \end{cases} 
        \]
        and
        \[
            S^\ell = \begin{cases}
                S^{\ell -1} \cup \{(i_\ell, j_{\ell -1})\} \qquad \text{ if } (i_\ell - u_{\ell-1}) = 0  \\
                S^{\ell -1} \cup \{(i_\ell -1, i_\ell)\} \qquad \text{otherwise.}
            \end{cases}
        \]
        \item[\underline{\textbf{Case 3:}}] $(i_\ell - u_{\ell-1})$ is odd and $(v_{\ell-1} - j_\ell)$ is even.

        We first define
        \[
            Q^\ell = \begin{cases}
                (P^{\ell-1} \setminus \{(i_\ell, j_\ell)\}) &\cup \quad \{(u_{\ell-1}+1, u_{\ell-1}+2), \ldots, (i_\ell -2, i_\ell -1)\} \\
                &\cup \quad \{(v_{\ell-1}-2, v_{\ell-1}-1), \ldots, (j_\ell, j_\ell +1)\} 
                \qquad \quad \text{ if } (v_{\ell-1} - j_\ell) > 0, \\
                (P^{\ell-1} \setminus \{(i_\ell, j_\ell)\}) &\cup \quad \{(u_{\ell-1}+1, u_{\ell-1}+2), \ldots, (i_\ell -2, i_\ell -1)\} \qquad \text{if } (v_{\ell-1} - j_\ell) = 0,
            \end{cases}
        \]
        
        and then define
        \[
            P^\ell = \begin{cases}
                Q^\ell \cup \{(i_\ell, i_\ell +1)\} \quad \text{ if } \quad i_\ell +1 \in R_-.\\
                Q^\ell \cup \{(i_\ell, j_\ell -1)\} \quad \text{ otherwise.} 
                \end{cases} 
        \]
        
        Also,
        \[
            u_\ell = \begin{cases}
                i_\ell +1 \quad \text{ if } \quad i_\ell +1 \in R_-\\
                i_\ell \qquad \quad \text{otherwise.} 
            \end{cases} 
            \qquad \qquad 
            v_\ell = \begin{cases}
                j_\ell \qquad \quad \text{if } \quad i_\ell +1 \in R_-\\
                j_\ell -1 \quad \text{ otherwise.} 
            \end{cases} 
        \]
        and
        \[
            S^\ell = \begin{cases}
                S^{\ell -1} \cup \{(i_{\ell -1}, j_\ell)\} \qquad \text{ if } (v_{\ell -1} - j_\ell) = 0.  \\
                S^{\ell -1} \cup \{(j_\ell, j_\ell +1)\} \qquad \text{otherwise.}
            \end{cases}
        \]

        \item[\underline{\textbf{Case 4:}}]  $(i_\ell - u_{\ell-1})$ and $(v_{\ell-1} - j_\ell)$ are both even.

        We first define
        \[
            Q^\ell = \begin{cases}
                (P^{\ell-1} \setminus \{(i_\ell, j_\ell)\}) \qquad \qquad \qquad \qquad  \qquad \qquad \qquad \qquad \text{ if } (i_\ell - u_{\ell-1}) = 0 \\
                (P^{\ell-1} \setminus \{(i_\ell, j_\ell)\}) \cup \{(u_{\ell-1}+1, u_{\ell-1}+2), \ldots, (i_\ell -1, i_\ell)\} \quad \text{ otherwise}
            \end{cases} 
        \]
        and then define
        \[
            P^\ell = \begin{cases}
                Q^\ell \qquad \qquad \qquad \qquad \qquad \qquad \qquad \qquad \text{ if } (v_{\ell-1} - j_\ell) = 0 \\
                Q^\ell \cup \{(v_{\ell-1}-2, v_{\ell-1}-1), \ldots, (j_\ell, j_\ell +1)\} \quad \text{ otherwise.}
            \end{cases}
        \]
        Also,
        \[
            u_\ell = i_\ell \qquad \qquad v_\ell = j_\ell \qquad \text{ and } \qquad S^\ell = \begin{cases}
                S^{\ell -1} \cup \{(i_\ell, j_{\ell -1})\} \qquad \text{ if } (i_\ell - u_{\ell-1}) = 0.  \\
                S^{\ell -1} \cup \{(i_\ell -1, i_\ell)\} \qquad \text{otherwise.}
            \end{cases}
        \]
    \end{enumerate}  
    Note that $(v_p - u_p)$ must be odd.
    So finally, we define $P = P^p \cup \{(u_p+1, u_p+2), \ldots, (v_p-2, v_p-1)\}$ and $S = S^p$.
    Firstly note that $P$ corresponds to an arc-partition, since $P^0$ is a partial arc-partition and it is easy to check that for each $\ell \in [p]$, $P^\ell$ is a valid extension of $P^{\ell-1}$.
    
    Since $p \geq d/8$, $S$ clearly has size at least $d/8$.
    The only thing left to complete the proof is to check that $S$ has the other required property.
    For any $\ell \in [p]$, let us assume that $S^{\ell -1} \subset P$ has the property that if $(i,j) \in S^{\ell -1}$, $w_i + w_j = 0$.    
    In Case 1, clearly, $S^\ell$ continues to have this property.
    In Cases 2 and 4, if $(i_\ell - u_{\ell-1}) > 0$, then again $S^\ell$ clearly continues to have this property.
    When $(i_\ell - u_{\ell-1}) = 0$, it must mean that $i_\ell = u_{\ell-1} = i_{\ell -1} + 1$ and so $i_{\ell -1} + 1 \not\in R_-$.
    That is, in the previous iteration Case 2 had been true and $(i_\ell, j_{\ell -1}) = (i_{\ell -1} + 1, j_{\ell -1})$ had been added to $P^{\ell -1}$.
    Therefore $S^\ell \subset P$ and, also, $w_{i_\ell} + w_{j_{\ell -1}} = 0$.
    A similar argument for Case 3 then completes the proof.
\end{proof}

\newcommand{\etalchar}[1]{$^{#1}$}
\begin{thebibliography}{DMPY12}

\bibitem[AFS{\etalchar{+}}18]{AndersonFSSV18}
Matthew Anderson, Michael~A. Forbes, Ramprasad Saptharishi, Amir Shpilka, and
  Ben~Lee Volk.
\newblock Identity testing and lower bounds for read-\emph{k} oblivious
  algebraic branching programs.
\newblock {\em {ACM} Trans. Comput. Theory}, 10(1):3:1--3:30, 2018.

\bibitem[AGKS15]{AgrawalGKS15}
Manindra Agrawal, Rohit Gurjar, Arpita Korwar, and Nitin Saxena.
\newblock Hitting-sets for {ROABP} and sum of set-multilinear circuits.
\newblock {\em {SIAM} J. Comput.}, 44(3):669--697, 2015.

\bibitem[AR16]{ArvindR16}
Vikraman Arvind and S.~Raja.
\newblock Some lower bound results for set-multilinear arithmetic computations.
\newblock {\em Chic. J. Theor. Comput. Sci.}, 2016, 2016.

\bibitem[BDS22]{mfcs/BhargavDS22}
C.~S. Bhargav, Sagnik Dutta, and Nitin Saxena.
\newblock Improved lower bound, and proof barrier, for constant depth algebraic
  circuits.
\newblock In Stefan Szeider, Robert Ganian, and Alexandra Silva, editors, {\em
  47th International Symposium on Mathematical Foundations of Computer Science,
  {MFCS} 2022, August 22-26, 2022, Vienna, Austria}, volume 241 of {\em
  LIPIcs}, pages 18:1--18:16. Schloss Dagstuhl - Leibniz-Zentrum f{\"{u}}r
  Informatik, 2022.

\bibitem[BDS24]{BDS23}
C.S. Bhargav, Prateek Dwivedi, and Nitin Saxena.
\newblock Lower bounds for the sum of small-size algebraic branching programs.
\newblock {\em To appear in the proceedings of the Annual Conference on Theory
  and Applications of Models of Computation}, 2024.
\newblock \url{https://www.cse.iitk.ac.in/users/nitin/papers/sumRO.pdf}.

\bibitem[BG22]{BhargavaG22}
Vishwas Bhargava and Sumanta Ghosh.
\newblock Improved hitting set for orbit of roabps.
\newblock {\em Comput. Complex.}, 31(2):15, 2022.

\bibitem[BS21]{BishtS21}
Pranav Bisht and Nitin Saxena.
\newblock Blackbox identity testing for sum of special roabps and its border
  class.
\newblock {\em Comput. Complex.}, 30(1):8, 2021.

\bibitem[B{\"{u}}r00]{burg}
Peter B{\"{u}}rgisser.
\newblock Cook's versus valiant's hypothesis.
\newblock {\em Theor. Comput. Sci.}, 235(1):71--88, 2000.

\bibitem[CKSV22]{ChatterjeeKSV22}
Prerona Chatterjee, Mrinal Kumar, Adrian She, and Ben~Lee Volk.
\newblock Quadratic lower bounds for algebraic branching programs and formulas.
\newblock {\em Comput. Complex.}, 31(2):8, 2022.

\bibitem[DMPY12]{DvirMPY12}
Zeev Dvir, Guillaume Malod, Sylvain Perifel, and Amir Yehudayoff.
\newblock Separating multilinear branching programs and formulas.
\newblock In Howard~J. Karloff and Toniann Pitassi, editors, {\em Proceedings
  of the 44th Symposium on Theory of Computing Conference, {STOC} 2012, New
  York, NY, USA, May 19 - 22, 2012}, pages 615--624. {ACM}, 2012.

\bibitem[FS13]{ForbesS13}
Michael~A. Forbes and Amir Shpilka.
\newblock Quasipolynomial-time identity testing of non-commutative and
  read-once oblivious algebraic branching programs.
\newblock In {\em 54th Annual {IEEE} Symposium on Foundations of Computer
  Science, {FOCS} 2013, 26-29 October, 2013, Berkeley, CA, {USA}}, pages
  243--252. {IEEE} Computer Society, 2013.

\bibitem[FSS14]{ForbesSS14}
Michael~A. Forbes, Ramprasad Saptharishi, and Amir Shpilka.
\newblock Hitting sets for multilinear read-once algebraic branching programs,
  in any order.
\newblock In David~B. Shmoys, editor, {\em Symposium on Theory of Computing,
  {STOC} 2014, New York, NY, USA, May 31 - June 03, 2014}, pages 867--875.
  {ACM}, 2014.

\bibitem[GG20]{guo2020improved}
Zeyu Guo and Rohit Gurjar.
\newblock Improved explicit hitting-sets for roabps.
\newblock In {\em Approximation, Randomization, and Combinatorial Optimization.
  Algorithms and Techniques (APPROX/RANDOM 2020)}. Schloss
  Dagstuhl-Leibniz-Zentrum f{\"u}r Informatik, 2020.

\bibitem[GKS17]{GurjarKS17}
Rohit Gurjar, Arpita Korwar, and Nitin Saxena.
\newblock Identity testing for constant-width, and any-order, read-once
  oblivious arithmetic branching programs.
\newblock {\em Theory Comput.}, 13(1):1--21, 2017.

\bibitem[GKST17]{GurjarKST17}
Rohit Gurjar, Arpita Korwar, Nitin Saxena, and Thomas Thierauf.
\newblock Deterministic identity testing for sum of read-once oblivious
  arithmetic branching programs.
\newblock {\em Comput. Complex.}, 26(4):835--880, 2017.

\bibitem[GR21]{GR21}
Purnata Ghosal and B.~V.~Raghavendra Rao.
\newblock Limitations of sums of bounded read formulas and abps.
\newblock In Rahul Santhanam and Daniil Musatov, editors, {\em Computer Science
  - Theory and Applications - 16th International Computer Science Symposium in
  Russia, {CSR} 2021, Sochi, Russia, June 28 - July 2, 2021, Proceedings},
  volume 12730 of {\em Lecture Notes in Computer Science}, pages 147--169.
  Springer, 2021.

\bibitem[GV20]{GurjarV20}
Rohit Gurjar and Ben~Lee Volk.
\newblock Pseudorandom bits for oblivious branching programs.
\newblock {\em {ACM} Trans. Comput. Theory}, 12(2):8:1--8:12, 2020.

\bibitem[Jan08]{Jansen08}
Maurice~J. Jansen.
\newblock Lower bounds for syntactically multilinear algebraic branching
  programs.
\newblock In Edward Ochmanski and Jerzy Tyszkiewicz, editors, {\em Mathematical
  Foundations of Computer Science 2008, 33rd International Symposium, {MFCS}
  2008, Torun, Poland, August 25-29, 2008, Proceedings}, volume 5162 of {\em
  Lecture Notes in Computer Science}, pages 407--418. Springer, 2008.

\bibitem[Kay12]{Kayal12}
Neeraj Kayal.
\newblock An exponential lower bound for the sum of powers of bounded degree
  polynomials.
\newblock {\em Electron. Colloquium Comput. Complex.}, {TR12-081}, 2012.

\bibitem[KLSS17]{KayalLSS17}
Neeraj Kayal, Nutan Limaye, Chandan Saha, and Srikanth Srinivasan.
\newblock An exponential lower bound for homogeneous depth four arithmetic
  formulas.
\newblock {\em {SIAM} J. Comput.}, 46(1):307--335, 2017.

\bibitem[KNS20]{KayalNS20}
Neeraj Kayal, Vineet Nair, and Chandan Saha.
\newblock Separation between read-once oblivious algebraic branching programs
  (roabps) and multilinear depth-three circuits.
\newblock {\em {ACM} Trans. Comput. Theory}, 12(1):2:1--2:27, 2020.

\bibitem[KS17]{/KumarS17}
Mrinal Kumar and Shubhangi Saraf.
\newblock On the power of homogeneous depth 4 arithmetic circuits.
\newblock {\em {SIAM} J. Comput.}, 46(1):336--387, 2017.

\bibitem[KS22]{KushS22}
Deepanshu Kush and Shubhangi Saraf.
\newblock Improved low-depth set-multilinear circuit lower bounds.
\newblock In Shachar Lovett, editor, {\em 37th Computational Complexity
  Conference, {CCC} 2022, July 20-23, 2022, Philadelphia, PA, {USA}}, volume
  234 of {\em LIPIcs}, pages 38:1--38:16. Schloss Dagstuhl - Leibniz-Zentrum
  f{\"{u}}r Informatik, 2022.

\bibitem[KS23]{KushS23}
Deepanshu Kush and Shubhangi Saraf.
\newblock Near-optimal set-multilinear formula lower bounds.
\newblock In Amnon Ta{-}Shma, editor, {\em 38th Computational Complexity
  Conference, {CCC} 2023, July 17-20, 2023, Warwick, {UK}}, volume 264 of {\em
  LIPIcs}, pages 15:1--15:33. Schloss Dagstuhl - Leibniz-Zentrum f{\"{u}}r
  Informatik, 2023.

\bibitem[KSS14]{KayalSS14}
Neeraj Kayal, Chandan Saha, and Ramprasad Saptharishi.
\newblock A super-polynomial lower bound for regular arithmetic formulas.
\newblock In David~B. Shmoys, editor, {\em Symposium on Theory of Computing,
  {STOC} 2014, New York, NY, USA, May 31 - June 03, 2014}, pages 146--153.
  {ACM}, 2014.

\bibitem[KST16]{KayalST16}
Neeraj Kayal, Chandan Saha, and S{\'{e}}bastien Tavenas.
\newblock An almost cubic lower bound for depth three arithmetic circuits.
\newblock In Ioannis Chatzigiannakis, Michael Mitzenmacher, Yuval Rabani, and
  Davide Sangiorgi, editors, {\em 43rd International Colloquium on Automata,
  Languages, and Programming, {ICALP} 2016, July 11-15, 2016, Rome, Italy},
  volume~55 of {\em LIPIcs}, pages 33:1--33:15. Schloss Dagstuhl -
  Leibniz-Zentrum f{\"{u}}r Informatik, 2016.

\bibitem[LST21]{LST1}
Nutan Limaye, Srikanth Srinivasan, and S{\'{e}}bastien Tavenas.
\newblock Superpolynomial lower bounds against low-depth algebraic circuits.
\newblock In {\em 62nd {IEEE} Annual Symposium on Foundations of Computer
  Science, {FOCS} 2021, Denver, CO, USA, February 7-10, 2022}, pages 804--814.
  {IEEE}, 2021.

\bibitem[LST22]{LST3}
Nutan Limaye, Srikanth Srinivasan, and S{\'{e}}bastien Tavenas.
\newblock On the partial derivative method applied to lopsided set-multilinear
  polynomials.
\newblock In Shachar Lovett, editor, {\em 37th Computational Complexity
  Conference, {CCC} 2022, July 20-23, 2022, Philadelphia, PA, {USA}}, volume
  234 of {\em LIPIcs}, pages 32:1--32:23. Schloss Dagstuhl - Leibniz-Zentrum
  f{\"{u}}r Informatik, 2022.

\bibitem[MU05]{mitzenmacher_upfal_2005}
Michael Mitzenmacher and Eli Upfal.
\newblock {\em Probability and Computing: Randomized Algorithms and
  Probabilistic Analysis}.
\newblock Cambridge University Press, 2005.

\bibitem[Nis91]{Nisan91}
Noam Nisan.
\newblock Lower bounds for non-commutative computation (extended abstract).
\newblock In Cris Koutsougeras and Jeffrey~Scott Vitter, editors, {\em
  Proceedings of the 23rd Annual {ACM} Symposium on Theory of Computing, May
  5-8, 1991, New Orleans, Louisiana, {USA}}, pages 410--418. {ACM}, 1991.

\bibitem[NW97]{NisanW97}
Noam Nisan and Avi Wigderson.
\newblock Lower bounds on arithmetic circuits via partial derivatives.
\newblock {\em Comput. Complex.}, 6(3):217--234, 1997.

\bibitem[Raz06]{Raz06}
Ran Raz.
\newblock Separation of multilinear circuit and formula size.
\newblock {\em Theory Comput.}, 2(6):121--135, 2006.

\bibitem[Raz13]{Raz-Tensor}
Ran Raz.
\newblock Tensor-rank and lower bounds for arithmetic formulas.
\newblock {\em J. {ACM}}, 60(6):40:1--40:15, 2013.

\bibitem[RR20]{RR20}
C.~Ramya and B.~V.~Raghavendra Rao.
\newblock Lower bounds for special cases of syntactic multilinear abps.
\newblock {\em Theor. Comput. Sci.}, 809:1--20, 2020.

\bibitem[RS05]{RazS05}
Ran Raz and Amir Shpilka.
\newblock Deterministic polynomial identity testing in non-commutative models.
\newblock {\em Comput. Complex.}, 14(1):1--19, 2005.

\bibitem[RY08]{RazY08}
Ran Raz and Amir Yehudayoff.
\newblock Balancing syntactically multilinear arithmetic circuits.
\newblock {\em Comput. Complex.}, 17(4):515--535, 2008.

\bibitem[RY09]{RazY09}
Ran Raz and Amir Yehudayoff.
\newblock Lower bounds and separations for constant depth multilinear circuits.
\newblock {\em Comput. Complex.}, 18(2):171--207, 2009.

\bibitem[Sap15]{Saptarishi-survey}
Ramprasad Saptharishi.
\newblock A survey of lower bounds in arithmetic circuit complexity.
\newblock {\em Github Survey}, 2015.

\bibitem[SY10]{ShpilkaY10}
Amir Shpilka and Amir Yehudayoff.
\newblock Arithmetic circuits: {A} survey of recent results and open questions.
\newblock {\em Found. Trends Theor. Comput. Sci.}, 5(3-4):207--388, 2010.

\bibitem[TLS22]{LST2}
S{\'{e}}bastien Tavenas, Nutan Limaye, and Srikanth Srinivasan.
\newblock Set-multilinear and non-commutative formula lower bounds for iterated
  matrix multiplication.
\newblock In Stefano Leonardi and Anupam Gupta, editors, {\em {STOC} '22: 54th
  Annual {ACM} {SIGACT} Symposium on Theory of Computing, Rome, Italy, June 20
  - 24, 2022}, pages 416--425. {ACM}, 2022.

\bibitem[Val79]{Val79}
L.~G. Valiant.
\newblock Completeness classes in algebra.
\newblock In {\em Proceedings of the Eleventh Annual ACM Symposium on Theory of
  Computing}, STOC '79, page 249–261, New York, NY, USA, 1979. Association
  for Computing Machinery.

\end{thebibliography}
\end{document}
